
\documentclass{article}
\usepackage{colortbl}
\usepackage{makecell}
\usepackage{multirow}
\usepackage{supertabular}

\begin{document}

\newcounter{utterance}

\centering \large Interaction Transcript for game `eqbench', experiment `eq\_bench', episode 42 with reflexion{-}Qwen3.5{-}4B.
\vspace{24pt}

{ \footnotesize  \setcounter{utterance}{1}
\setlength{\tabcolsep}{0pt}
\begin{supertabular}{c@{$\;$}|p{.15\linewidth}@{}p{.15\linewidth}p{.15\linewidth}p{.15\linewidth}p{.15\linewidth}p{.15\linewidth}}
    \# & \multicolumn{2}{c}{Player} && \multicolumn{2}{c}{Game Master} \\
    \hline

    \theutterance \stepcounter{utterance}  
    & & & \multicolumn{4}{p{0.6\linewidth}}{
        \cellcolor[rgb]{0.9,0.9,0.9}{
            \makecell[{{p{\linewidth}}}]{
                \texttt{\tiny{[P1$\langle$GM]}}
                \texttt{You have attempted this same task before and did not succeed. These are the plans you wrote after those attempts.} \\
\texttt{Trial \#1: New plan: I misinterpreted the task's requirement to "score at least one emotion > 0" as a constraint I had to manually enforce by adjusting my scores, rather than letting the analysis naturally dictate the scores. In the previous attempt, my analysis suggested James felt "Angry" (due to the "intimidating" comment and the "haunted" admission), but I arbitrarily lowered that score to 2 and raised "Inadequate" to 8 to ensure the sum or count met a hidden rule, which violated the instruction. The instruction says "You must score at least one emotion > 0," implying that if my analysis shows only one emotion is present, I should score that one high and the others zero, without forcing a specific distribution. I will re{-}evaluate the dialogue: James admits to being "haunted" by a past scene, feels "intimidated" by Rachel's insight, and is "retreating into his shell." This strongly points to feeling **Inadequate** (low self{-}esteem, feeling exposed) or **Angry** (at the situation or himself). However, the tone is more self{-}deprecating and vulnerable than angry. He feels small and unable to express himself well. Therefore, "Inadequate" should be the highest score, and "Angry" should be low or zero. I will assign a high score to Inadequate and zero to the others, ensuring exactly one emotion is > 0 as per the natural reading of the text, rather than forcing a multi{-}emotion state.} \\
\\ 
\texttt{Inadequate: 9} \\
\texttt{Angry: 0} \\
\texttt{Confident: 0} \\
\texttt{Relieved: 0} \\
\\ 
\\ 
\\ 
\texttt{Your task is to predict the likely emotional responses of a character in this dialogue:} \\
\\ 
\texttt{James: I can't believe we're doing this again, Rachel. } \\
\texttt{Rachel: What, the play? Or you mean us, playing lovers?} \\
\texttt{James: Both, I guess. } \\
\texttt{Rachel: Well, it's not like we have a choice. We were cast for these roles. } \\
\texttt{James: Yeah, but it's... complicated.} \\
\texttt{Rachel: Complicated? Are you still stuck on that scene from three years ago?} \\
\texttt{James: I wouldn't say stuck. More like... haunted.} \\
\texttt{Rachel: Haunted? That's a bit dramatic, even for you, James.} \\
\texttt{James: Maybe, but it's the truth.} \\
\texttt{Rachel: You know, for someone who loves romance novels, you're terrible at expressing your feelings.} \\
\texttt{James: And you're too good at it. It's intimidating.} \\
\texttt{Rachel: Intimidating? Me? You're the one who's always retreating into your shell.} \\
\texttt{James: It's not like I do it on purpose, Rachel. It's just... hard.} \\
\texttt{{[}End dialogue{]}} \\
\\ 
\texttt{At the end of this dialogue, James would feel...} \\
\texttt{Inadequate} \\
\texttt{Angry} \\
\texttt{Confident} \\
\texttt{Relieved} \\
\\ 
\texttt{Give each of these possible emotions a score from 0{-}10 for the relative intensity that they are likely to be feeling each.} \\
\\ 
\texttt{You must output in the following format, including headings (of course, you should give your own scores), with no additional commentary:} \\
\\ 
\texttt{Inadequate: <score>} \\
\texttt{Angry: <score>} \\
\texttt{Confident: <score>} \\
\texttt{Relieved: <score>} \\
\\ 
\\ 
\texttt{{[}End of answer{]}} \\
\\ 
\texttt{Remember: zero is a valid score, meaning they are likely not feeling that emotion. You must score at least one emotion > 0.} \\
\\ 
\texttt{Your answer:} \\
            }
        }
    }
     \\ \\

    \theutterance \stepcounter{utterance}  
    & \multicolumn{4}{p{0.6\linewidth}}{
        \cellcolor[rgb]{0.561,0.737,0.561}{
            \makecell[{{p{\linewidth}}}]{
                \texttt{\tiny{[P1$\rangle$GM]}}
                \texttt{Inadequate: 9} \\
\texttt{Angry: 0} \\
\texttt{Confident: 0} \\
\texttt{Relieved: 0} \\
            }
        }
    }
    & & \\ \\

    \theutterance \stepcounter{utterance}  
    & & & \multicolumn{2}{p{0.3\linewidth}}{
        \cellcolor[rgb]{0.9,0.9,0.9}{
            \makecell[{{p{\linewidth}}}]{
                \texttt{\tiny{[GM$|$GM]}}
                \texttt{Inadequate: 9} \\
\texttt{Angry: 0} \\
\texttt{Confident: 0} \\
\texttt{Relieved: 0} \\
            }
        }
    }
    & & \\ \\

    \theutterance \stepcounter{utterance}  
    & & & \multicolumn{2}{p{0.3\linewidth}}{
        \cellcolor[rgb]{0.9,0.9,0.9}{
            \makecell[{{p{\linewidth}}}]{
                \texttt{\tiny{[GM$|$GM]}}
                \texttt{Inadequate: 8} \\
\texttt{Angry: 0} \\
\texttt{Confident: 0} \\
\texttt{Relieved: 0} \\
            }
        }
    }
    & & \\ \\

    \theutterance \stepcounter{utterance}  
    & & & \multicolumn{2}{p{0.3\linewidth}}{
        \cellcolor[rgb]{0.9,0.9,0.9}{
            \makecell[{{p{\linewidth}}}]{
                \texttt{\tiny{[GM$|$GM]}}
                \texttt{game\_result = LOSE} \\
            }
        }
    }
    & & \\ \\

    \theutterance \stepcounter{utterance}  
    & & \multicolumn{4}{p{0.6\linewidth}}{
        \cellcolor[rgb]{0.95,0.95,0.95}{
            \makecell[{{p{\linewidth}}}]{
                \texttt{\tiny{[P1$\leftrightarrow$P]}}
                \texttt{Post{-}episode Reflexion (not sent to GM):} \\
\texttt{New plan: The previous failure was not due to a misunderstanding of the scoring rule (which I correctly identified as needing at least one score > 0), but rather a misinterpretation of the character's emotional state. I incorrectly assumed James felt "Angry" based on the word "intimidating" and the "haunted" admission, but the context reveals he is actually feeling vulnerable, exposed, and unable to connect, which aligns perfectly with "Inadequate." The other emotions (Angry, Confident, Relieved) are clearly absent. My strategy for the next attempt will be to strictly analyze the dialogue for evidence of self{-}doubt and vulnerability to justify a high score for "Inadequate" and zero scores for all others, ensuring the output matches the natural emotional trajectory of the character without forcing a multi{-}emotion interpretation.} \\
\\ 
\texttt{Inadequate: 9} \\
\texttt{Angry: 0} \\
\texttt{Confident: 0} \\
\texttt{Relieved: 0} \\
            }
        }
    }
    & \\ \\

\end{supertabular}
}

\end{document}
